% ProveIt Technical Manual
% Compiled with XeLaTeX for advanced typography and Unicode support
%
% Build: xelatex proveit-manual.tex && xelatex proveit-manual.tex
% (Run twice for cross-references and index)

\documentclass[11pt,a4paper,twoside]{report}

%===============================================================================
% XETEX-SPECIFIC PACKAGES
%===============================================================================
\usepackage{fontspec}
\usepackage{unicode-math}
\usepackage{polyglossia}
\setdefaultlanguage{english}

% Typography - Use system fonts with OpenType features
\setmainfont{Latin Modern Roman}[
    Ligatures=TeX,
    Numbers=OldStyle
]
\setsansfont{Latin Modern Sans}[
    Ligatures=TeX
]
\setmonofont{Latin Modern Mono}[
    Scale=0.9
]
\setmathfont{Latin Modern Math}

%===============================================================================
% ESSENTIAL PACKAGES
%===============================================================================
\usepackage{geometry}
\geometry{
    paper=a4paper,
    inner=25mm,
    outer=20mm,
    top=25mm,
    bottom=25mm,
    bindingoffset=10mm
}

\usepackage{graphicx}
\usepackage{xcolor}
\usepackage{hyperref}
\usepackage{cleveref}
\usepackage{booktabs}
\usepackage{longtable}
\usepackage{multirow}
\usepackage{enumitem}
\usepackage{fancyhdr}
\usepackage{titlesec}
\usepackage{tocloft}
\usepackage{makeidx}
\usepackage{imakeidx}

%===============================================================================
% CODE LISTINGS
%===============================================================================
\usepackage{listings}
\usepackage{lstfiracode}

\definecolor{codebg}{HTML}{F8F8F8}
\definecolor{codeframe}{HTML}{E0E0E0}
\definecolor{keyword}{HTML}{0000FF}
\definecolor{comment}{HTML}{008000}
\definecolor{string}{HTML}{A31515}
\definecolor{identifier}{HTML}{001080}

\lstdefinelanguage{Rust}{
    keywords={as, break, const, continue, crate, else, enum, extern, false, fn, for, if, impl, in, let, loop, match, mod, move, mut, pub, ref, return, self, Self, static, struct, super, trait, true, type, unsafe, use, where, while, async, await, dyn, abstract, become, box, do, final, macro, override, priv, typeof, unsized, virtual, yield},
    keywordstyle=\color{keyword}\bfseries,
    ndkeywords={bool, char, i8, i16, i32, i64, i128, isize, u8, u16, u32, u64, u128, usize, f32, f64, str, String, Vec, Option, Result, Some, None, Ok, Err, Box, Rc, Arc, Cell, RefCell, HashMap, HashSet, BTreeMap, BTreeSet},
    ndkeywordstyle=\color{identifier}\bfseries,
    identifierstyle=\color{black},
    sensitive=true,
    comment=[l]{//},
    morecomment=[s]{/*}{*/},
    commentstyle=\color{comment}\itshape,
    stringstyle=\color{string},
    morestring=[b]",
    morestring=[b]',
}

\lstset{
    language=Rust,
    basicstyle=\ttfamily\small,
    backgroundcolor=\color{codebg},
    frame=single,
    rulecolor=\color{codeframe},
    framesep=5pt,
    breaklines=true,
    breakatwhitespace=true,
    tabsize=4,
    showstringspaces=false,
    numbers=left,
    numberstyle=\tiny\color{gray},
    numbersep=10pt,
    xleftmargin=15pt,
    xrightmargin=5pt,
    aboveskip=10pt,
    belowskip=10pt,
    captionpos=b,
    escapeinside={(*@}{@*)},
}

%===============================================================================
% TIKZ FOR DIAGRAMS
%===============================================================================
\usepackage{tikz}
\usetikzlibrary{
    arrows.meta,
    shapes.geometric,
    shapes.symbols,
    positioning,
    calc,
    fit,
    backgrounds,
    decorations.pathmorphing,
    decorations.markings,
    matrix,
    chains,
    scopes,
    cd  % Commutative diagrams
}

% Custom TikZ styles for ProveIt diagrams
\tikzstyle{module} = [
    rectangle,
    rounded corners=3pt,
    draw=blue!50,
    fill=blue!10,
    minimum width=2.5cm,
    minimum height=1cm,
    text centered,
    font=\sffamily\small
]

\tikzstyle{trait} = [
    rectangle,
    draw=green!50!black,
    fill=green!10,
    minimum width=2cm,
    minimum height=0.8cm,
    text centered,
    font=\sffamily\small\itshape
]

\tikzstyle{proofnode} = [
    circle,
    draw=purple!50,
    fill=purple!10,
    minimum size=0.8cm,
    font=\small
]

\tikzstyle{arrow} = [
    ->,
    >=Stealth,
    thick
]

%===============================================================================
% THEOREM ENVIRONMENTS
%===============================================================================
\usepackage{amsthm}
\usepackage{thmtools}

\definecolor{theorembg}{HTML}{F0F7FF}
\definecolor{theoremframe}{HTML}{4A90D9}
\definecolor{definitionbg}{HTML}{F0FFF0}
\definecolor{definitionframe}{HTML}{4AD94A}
\definecolor{examplebg}{HTML}{FFF8F0}
\definecolor{exampleframe}{HTML}{D9944A}

\declaretheoremstyle[
    headfont=\bfseries\sffamily\color{theoremframe},
    bodyfont=\normalfont,
    spaceabove=12pt,
    spacebelow=12pt,
    mdframed={
        backgroundcolor=theorembg,
        linecolor=theoremframe,
        linewidth=2pt,
        leftline=true,
        rightline=false,
        topline=false,
        bottomline=false,
        innerleftmargin=10pt,
        innerrightmargin=10pt,
        innertopmargin=10pt,
        innerbottommargin=10pt
    }
]{theoremstyle}

\declaretheoremstyle[
    headfont=\bfseries\sffamily\color{definitionframe},
    bodyfont=\normalfont,
    spaceabove=12pt,
    spacebelow=12pt,
    mdframed={
        backgroundcolor=definitionbg,
        linecolor=definitionframe,
        linewidth=2pt,
        leftline=true,
        rightline=false,
        topline=false,
        bottomline=false,
        innerleftmargin=10pt,
        innerrightmargin=10pt,
        innertopmargin=10pt,
        innerbottommargin=10pt
    }
]{definitionstyle}

\declaretheoremstyle[
    headfont=\bfseries\sffamily\color{exampleframe},
    bodyfont=\normalfont,
    spaceabove=12pt,
    spacebelow=12pt,
    mdframed={
        backgroundcolor=examplebg,
        linecolor=exampleframe,
        linewidth=2pt,
        leftline=true,
        rightline=false,
        topline=false,
        bottomline=false,
        innerleftmargin=10pt,
        innerrightmargin=10pt,
        innertopmargin=10pt,
        innerbottommargin=10pt
    }
]{examplestyle}

\usepackage{mdframed}

\declaretheorem[style=theoremstyle, name=Theorem, numberwithin=chapter]{theorem}
\declaretheorem[style=theoremstyle, name=Lemma, sibling=theorem]{lemma}
\declaretheorem[style=theoremstyle, name=Proposition, sibling=theorem]{proposition}
\declaretheorem[style=theoremstyle, name=Corollary, sibling=theorem]{corollary}
\declaretheorem[style=definitionstyle, name=Definition, sibling=theorem]{definition}
\declaretheorem[style=examplestyle, name=Example, sibling=theorem]{example}

%===============================================================================
% MATHEMATICAL NOTATION
%===============================================================================
% Type theory notation
\newcommand{\Type}{\mathsf{Type}}
\newcommand{\Prop}{\mathsf{Prop}}
\newcommand{\Set}{\mathsf{Set}}
\newcommand{\Unit}{\mathbf{1}}
\newcommand{\Empty}{\mathbf{0}}
\newcommand{\Bool}{\mathsf{Bool}}
\newcommand{\Nat}{\mathbb{N}}
\newcommand{\Int}{\mathbb{Z}}
\newcommand{\Real}{\mathbb{R}}

% Dependent types
\newcommand{\Sigma}[2]{\sum_{#1} #2}
\newcommand{\Pi}[2]{\prod_{#1} #2}
\newcommand{\lmark}[1]{\lambda #1.\,}
\newcommand{\app}[2]{#1 \, #2}

% Identity types
\newcommand{\Id}[3]{\mathsf{Id}_{#1}(#2, #3)}
\newcommand{\refl}[1]{\mathsf{refl}_{#1}}
\newcommand{\transport}[3]{\mathsf{transport}^{#1}_{#2}(#3)}
\newcommand{\ap}[2]{\mathsf{ap}_{#1}(#2)}

% Homotopy type theory
\newcommand{\isEquiv}[1]{\mathsf{isEquiv}(#1)}
\newcommand{\Equiv}[2]{#1 \simeq #2}
\newcommand{\fib}[2]{\mathsf{fib}_{#1}(#2)}
\newcommand{\isContr}[1]{\mathsf{isContr}(#1)}
\newcommand{\isProp}[1]{\mathsf{isProp}(#1)}
\newcommand{\isSet}[1]{\mathsf{isSet}(#1)}
\newcommand{\trunc}[2]{\lVert #2 \rVert_{#1}}
\newcommand{\proptrunc}[1]{\trunc{-1}{#1}}

% Category theory
\newcommand{\Hom}[3]{\mathsf{Hom}_{#1}(#2, #3)}
\newcommand{\id}[1]{\mathsf{id}_{#1}}
\newcommand{\op}[1]{#1^{\mathsf{op}}}
\newcommand{\Cat}{\mathbf{Cat}}
\newcommand{\Func}[2]{\mathsf{Func}(#1, #2)}
\newcommand{\Nat}{\mathsf{Nat}}

% Proof notation
\newcommand{\proves}{\vdash}
\newcommand{\entails}{\models}
\newcommand{\inmark}[2]{#1 : #2}
\newcommand{\ctx}{\Gamma}
\newcommand{\subst}[3]{#1[#2 := #3]}

%===============================================================================
% HEADERS AND FOOTERS
%===============================================================================
\pagestyle{fancy}
\fancyhf{}
\fancyhead[LE]{\small\itshape\leftmark}
\fancyhead[RO]{\small\itshape\rightmark}
\fancyfoot[C]{\thepage}
\renewcommand{\headrulewidth}{0.4pt}
\renewcommand{\footrulewidth}{0pt}

%===============================================================================
% HYPERREF CONFIGURATION
%===============================================================================
\hypersetup{
    colorlinks=true,
    linkcolor=blue!70!black,
    citecolor=green!50!black,
    urlcolor=purple!70!black,
    pdftitle={ProveIt Technical Manual},
    pdfauthor={ProveIt Contributors},
    pdfsubject={Geometric Formal Verification},
    pdfkeywords={formal verification, type theory, category theory, homotopy type theory, accessibility}
}

%===============================================================================
% INDEX
%===============================================================================
\makeindex[columns=2, title=Index, intoc]

%===============================================================================
% DOCUMENT INFORMATION
%===============================================================================
\title{
    \vspace{-2cm}
    {\Huge\sffamily\bfseries ProveIt}\\[0.5cm]
    {\Large\sffamily Technical Manual}\\[1cm]
    {\large Geometric Construction Environment for\\Accessible Formal Verification}
}
\author{
    \textsc{ProveIt Contributors}\\[0.3cm]
    \small\url{https://github.com/TensorHusker/ProveIt}
}
\date{
    Version 0.1.0\\[0.2cm]
    \today
}

%===============================================================================
% DOCUMENT BODY
%===============================================================================
\begin{document}

\maketitle
\thispagestyle{empty}

\clearpage
\tableofcontents
\clearpage
\listoffigures
\listoftables
\clearpage

%-------------------------------------------------------------------------------
\chapter{Introduction}
\label{chap:introduction}
%-------------------------------------------------------------------------------

\section{Vision}
\index{vision}

ProveIt\index{ProveIt} reimagines formal verification through the lens of \emph{geometric construction}. Instead of wrestling with abstract syntax trees and type annotations, users build proofs by constructing geometric objects that embody logical relationships.

\begin{definition}[Geometric Proof Construction]
\label{def:geometric-proof}
A \emph{geometric proof construction} is a sequence of geometric operations $(g_1, g_2, \ldots, g_n)$ where each operation $g_i$ corresponds to a valid inference rule in the underlying formal system, and the spatial relationships between constructed objects encode logical dependencies.
\end{definition}

\subsection{Core Principles}

\begin{enumerate}[label=\textbf{\arabic*.}, leftmargin=*, itemsep=10pt]
    \item \textbf{Accessibility-First Design}\index{accessibility}

    All features must be accessible to blind, neurodivergent, and post-injury users. This is not an afterthought but a foundational design constraint.

    \item \textbf{Multiple Formal Methods}\index{formal methods}

    Support for type theory\index{type theory}, category theory\index{category theory}, and homotopy type theory\index{homotopy type theory}, with seamless translation between formalisms.

    \item \textbf{Real-Time Verification}\index{verification!real-time}

    Proofs are validated as they are constructed, providing immediate feedback on validity.

    \item \textbf{Spatial Reasoning}\index{spatial reasoning}

    Geometric constructions serve as the primary interaction paradigm, leveraging human spatial intuition.
\end{enumerate}

%-------------------------------------------------------------------------------
\chapter{Mathematical Foundations}
\label{chap:foundations}
%-------------------------------------------------------------------------------

\section{Type Theory}
\index{type theory}

ProveIt's foundation rests on dependent type theory\index{type theory!dependent}, which provides a unified framework for logic and computation.

\begin{definition}[Dependent Function Type]
\label{def:pi-type}
Given a type $A : \Type$ and a family $B : A \to \Type$, the \emph{dependent function type} (or $\Pi$-type) is:
\[
    \Pi_{x : A} B(x) \quad \text{or equivalently} \quad (x : A) \to B(x)
\]
Elements of this type are functions $f$ such that for any $a : A$, we have $f(a) : B(a)$.
\end{definition}

\begin{definition}[Dependent Pair Type]
\label{def:sigma-type}
Given a type $A : \Type$ and a family $B : A \to \Type$, the \emph{dependent pair type} (or $\Sigma$-type) is:
\[
    \Sigma_{x : A} B(x) \quad \text{or equivalently} \quad (x : A) \times B(x)
\]
Elements are pairs $(a, b)$ where $a : A$ and $b : B(a)$.
\end{definition}

\subsection{Type Judgments}
\index{type theory!judgments}

The fundamental judgments in our type system are:

\begin{table}[htbp]
\centering
\caption{Type-theoretic judgments in ProveIt}
\label{tab:judgments}
\begin{tabular}{@{}lll@{}}
\toprule
\textbf{Judgment} & \textbf{Reading} & \textbf{Meaning} \\
\midrule
$\ctx \proves A : \Type$ & ``$A$ is a type in context $\ctx$'' & Type formation \\
$\ctx \proves a : A$ & ``$a$ has type $A$ in context $\ctx$'' & Term introduction \\
$\ctx \proves a \equiv b : A$ & ``$a$ equals $b$ at type $A$'' & Definitional equality \\
$\ctx \proves \Id{A}{a}{b}$ & ``$a$ is propositionally equal to $b$'' & Identity type \\
\bottomrule
\end{tabular}
\end{table}

\subsection{Inference Rules as Geometric Operations}
\index{inference rules}

Each inference rule corresponds to a geometric construction:

\begin{figure}[htbp]
\centering
\begin{tikzpicture}[node distance=1.5cm]
    % Function application
    \node[proofnode] (f) {$f$};
    \node[proofnode, right=of f] (a) {$a$};
    \node[proofnode, below=of $(f)!0.5!(a)$] (fa) {$f(a)$};

    \draw[arrow] (f) -- (fa) node[midway, left, font=\scriptsize] {apply};
    \draw[arrow] (a) -- (fa);

    % Lambda abstraction
    \node[proofnode, right=3cm of a] (body) {$e$};
    \node[proofnode, below=of body] (lam) {$\lambda x.e$};
    \node[draw=none, above=0.5cm of body, font=\scriptsize] {$[x:A]$};

    \draw[arrow, dashed] (body) -- (lam) node[midway, right, font=\scriptsize] {abstract};

    % Labels
    \node[below=0.5cm of fa, font=\sffamily\small] {Application};
    \node[below=0.5cm of lam, font=\sffamily\small] {Abstraction};
\end{tikzpicture}
\caption{Geometric representation of $\lambda$-calculus operations}
\label{fig:lambda-geometry}
\end{figure}

\section{Homotopy Type Theory}
\index{homotopy type theory}

ProveIt implements homotopy type theory (HoTT)\index{HoTT|see{homotopy type theory}}, which interprets types as spaces and terms as points.

\begin{theorem}[Univalence Axiom]
\label{thm:univalence}
\index{univalence}
For types $A, B : \Type$, the canonical map
\[
    \mathsf{idtoeqv} : (A = B) \to (A \simeq B)
\]
is an equivalence.
\end{theorem}

The univalence axiom has profound implications for formal verification:

\begin{corollary}
Equivalent types are interchangeable in all contexts. If $\Equiv{A}{B}$ and $P : \Type \to \Type$, then $P(A) \simeq P(B)$.
\end{corollary}

\subsection{Higher Inductive Types}
\index{higher inductive types}

ProveIt supports higher inductive types (HITs), which allow constructors for paths:

\begin{example}[Circle as HIT]
\label{ex:circle}
The circle $S^1$ is defined as:
\begin{itemize}
    \item A point constructor: $\mathsf{base} : S^1$
    \item A path constructor: $\mathsf{loop} : \mathsf{base} = \mathsf{base}$
\end{itemize}
\end{example}

\begin{figure}[htbp]
\centering
\begin{tikzpicture}
    \draw[thick] (0,0) circle (1.5cm);
    \fill (1.5,0) circle (3pt) node[right] {$\mathsf{base}$};
    \draw[->, thick, blue, decorate, decoration={snake, amplitude=1pt, segment length=5pt}]
        (1.4,0.3) arc (15:345:1.5cm);
    \node[blue] at (0, 2) {$\mathsf{loop}$};
\end{tikzpicture}
\caption{The circle $S^1$ as a higher inductive type}
\label{fig:circle}
\end{figure}

\section{Category Theory}
\index{category theory}

ProveIt's internal language can be interpreted categorically, enabling powerful abstraction.

\begin{definition}[Category]
\label{def:category}
A \emph{category} $\mathcal{C}$ consists of:
\begin{itemize}
    \item A collection of \emph{objects} $\mathsf{Ob}(\mathcal{C})$
    \item For each pair of objects $A, B$, a set $\Hom{\mathcal{C}}{A}{B}$ of \emph{morphisms}
    \item Identity morphisms $\id{A} : \Hom{\mathcal{C}}{A}{A}$
    \item Composition $\circ : \Hom{\mathcal{C}}{B}{C} \times \Hom{\mathcal{C}}{A}{B} \to \Hom{\mathcal{C}}{A}{C}$
\end{itemize}
satisfying associativity and identity laws.
\end{definition}

\begin{figure}[htbp]
\centering
\begin{tikzcd}[row sep=large, column sep=large]
    A \arrow[r, "f"] \arrow[dr, "g \circ f"'] & B \arrow[d, "g"] \\
    & C
\end{tikzcd}
\caption{Composition of morphisms in a category}
\label{fig:composition}
\end{figure}

%-------------------------------------------------------------------------------
\chapter{Architecture}
\label{chap:architecture}
%-------------------------------------------------------------------------------

\section{Module Organization}
\index{architecture}

\begin{figure}[htbp]
\centering
\begin{tikzpicture}[node distance=1.2cm and 2cm]
    % Core modules
    \node[module] (lib) {lib.rs};

    \node[module, below left=of lib] (proof) {proof/};
    \node[module, below=of lib] (geometry) {geometry/};
    \node[module, below right=of lib] (formal) {formal\_methods/};

    \node[module, below=of proof] (access) {accessibility/};
    \node[module, below=of geometry] (verify) {verification/};

    % Traits
    \node[trait, right=2.5cm of formal] (verifiable) {Verifiable};
    \node[trait, below=0.8cm of verifiable] (construct) {Constructible};
    \node[trait, below=0.8cm of construct] (accessible) {Accessible};

    % Connections
    \draw[arrow] (lib) -- (proof);
    \draw[arrow] (lib) -- (geometry);
    \draw[arrow] (lib) -- (formal);
    \draw[arrow] (proof) -- (access);
    \draw[arrow] (geometry) -- (verify);
    \draw[arrow] (proof) -- (verify);

    \draw[arrow, dashed, blue!60] (formal) -- (verifiable);
    \draw[arrow, dashed, blue!60] (geometry) -- (construct);
    \draw[arrow, dashed, blue!60] (access) -- (accessible);

    % Legend
    \node[below=3cm of verify, font=\sffamily\small] {
        \tikz{\node[module, minimum width=1cm, minimum height=0.5cm] {};} Module \quad
        \tikz{\node[trait, minimum width=1cm, minimum height=0.5cm] {};} Trait
    };
\end{tikzpicture}
\caption{ProveIt module architecture}
\label{fig:architecture}
\end{figure}

\section{Core Traits}
\index{traits}

\begin{lstlisting}[caption={Core trait definitions}, label=lst:traits]
/// A type that can be formally verified
pub trait Verifiable {
    /// The type of proof witness
    type Witness;

    /// Verify the construction, returning a witness on success
    fn verify(&self) -> Result<Self::Witness, VerificationError>;

    /// Check if verification is currently valid
    fn is_valid(&self) -> bool;
}

/// A type that can be geometrically constructed
pub trait Constructible: Verifiable {
    /// The type of construction steps
    type Step;

    /// Apply a construction step
    fn apply_step(&mut self, step: Self::Step) -> Result<(), ConstructionError>;

    /// Get the construction history
    fn history(&self) -> &[Self::Step];
}

/// A type with accessibility metadata
pub trait Accessible {
    /// Get a screen-reader friendly description
    fn describe(&self) -> AccessibleDescription;

    /// Get audio cue for this element
    fn audio_cue(&self) -> Option<AudioCue>;

    /// Get haptic pattern for this element
    fn haptic_pattern(&self) -> Option<HapticPattern>;
}
\end{lstlisting}

\section{Design Patterns}
\index{design patterns}

\subsection{Builder Pattern for Proofs}
\index{design patterns!builder}

\begin{lstlisting}[caption={Proof builder pattern}, label=lst:builder]
pub struct ProofBuilder<S: ProofState> {
    context: Context,
    steps: Vec<ProofStep>,
    state: PhantomData<S>,
}

// Type states for proof construction
pub struct Unstarted;
pub struct InProgress;
pub struct Complete;

impl ProofBuilder<Unstarted> {
    pub fn new() -> Self { /* ... */ }

    pub fn assume(self, prop: Proposition) -> ProofBuilder<InProgress> {
        /* ... */
    }
}

impl ProofBuilder<InProgress> {
    pub fn apply_rule(self, rule: InferenceRule) -> Self { /* ... */ }

    pub fn finish(self) -> Result<ProofBuilder<Complete>, ProofError> {
        /* ... */
    }
}

impl ProofBuilder<Complete> {
    pub fn extract(self) -> Proof { /* ... */ }
}
\end{lstlisting}

%-------------------------------------------------------------------------------
\chapter{Accessibility}
\label{chap:accessibility}
%-------------------------------------------------------------------------------

\section{Design Philosophy}
\index{accessibility!design}

ProveIt's accessibility is not an add-on but a core architectural concern. Every feature is designed with multiple modalities in mind.

\begin{table}[htbp]
\centering
\caption{Accessibility modalities in ProveIt}
\label{tab:modalities}
\begin{tabular}{@{}p{3cm}p{4cm}p{5cm}@{}}
\toprule
\textbf{Modality} & \textbf{Use Case} & \textbf{Implementation} \\
\midrule
Visual & Default interaction & Standard GUI with high contrast options \\
Auditory & Screen reader users & ARIA labels, audio cues for operations \\
Haptic & Enhanced spatial awareness & Vibration patterns for geometric relationships \\
Keyboard & Motor accessibility & Full keyboard navigation, customizable shortcuts \\
\bottomrule
\end{tabular}
\end{table}

\section{WCAG Compliance}
\index{accessibility!WCAG}
\index{WCAG}

ProveIt targets WCAG 2.1 Level AAA compliance:

\begin{itemize}[itemsep=5pt]
    \item \textbf{Perceivable}: All information is presentable in multiple modalities
    \item \textbf{Operable}: All functionality is keyboard accessible
    \item \textbf{Understandable}: Clear, consistent interface with predictable behavior
    \item \textbf{Robust}: Compatible with assistive technologies
\end{itemize}

\subsection{Audio Representation of Proofs}
\index{accessibility!audio}

Geometric constructions are mapped to audio representations:

\begin{lstlisting}[caption={Audio cue generation}, label=lst:audio]
impl Accessible for GeometricPrimitive {
    fn audio_cue(&self) -> Option<AudioCue> {
        Some(match self {
            GeometricPrimitive::Point(p) => AudioCue::Tone {
                frequency: map_to_frequency(p.x),
                duration: Duration::from_millis(100),
                pan: map_to_stereo(p.y),
            },
            GeometricPrimitive::Line(l) => AudioCue::Sweep {
                start_freq: map_to_frequency(l.start.x),
                end_freq: map_to_frequency(l.end.x),
                duration: Duration::from_millis(300),
            },
            GeometricPrimitive::Circle(c) => AudioCue::Pulse {
                frequency: map_to_frequency(c.center.x),
                rate: map_to_rate(c.radius),
                duration: Duration::from_millis(500),
            },
        })
    }
}
\end{lstlisting}

%-------------------------------------------------------------------------------
\chapter{Code Style Guidelines}
\label{chap:style}
%-------------------------------------------------------------------------------

\section{Rust Conventions}
\index{code style}

\subsection{Naming}

\begin{table}[htbp]
\centering
\caption{Naming conventions}
\label{tab:naming}
\begin{tabular}{@{}lll@{}}
\toprule
\textbf{Item} & \textbf{Convention} & \textbf{Example} \\
\midrule
Types & PascalCase & \texttt{ProofStep}, \texttt{GeometricPrimitive} \\
Traits & PascalCase & \texttt{Verifiable}, \texttt{Constructible} \\
Functions & snake\_case & \texttt{verify\_construction} \\
Variables & snake\_case & \texttt{proof\_context} \\
Constants & SCREAMING\_SNAKE\_CASE & \texttt{MAX\_PROOF\_DEPTH} \\
Modules & snake\_case & \texttt{geometric\_primitives} \\
\bottomrule
\end{tabular}
\end{table}

\subsection{Error Handling}
\index{error handling}

\begin{lstlisting}[caption={Domain-specific error types}, label=lst:errors]
use thiserror::Error;

#[derive(Error, Debug)]
pub enum VerificationError {
    #[error("type mismatch: expected {expected}, found {found}")]
    TypeMismatch {
        expected: TypeId,
        found: TypeId,
        #[source]
        source: Option<Box<dyn std::error::Error + Send + Sync>>,
    },

    #[error("proof step {step} is invalid: {reason}")]
    InvalidStep {
        step: usize,
        reason: String,
    },

    #[error("construction incomplete: missing {missing:?}")]
    Incomplete {
        missing: Vec<String>,
    },
}
\end{lstlisting}

\section{Documentation Standards}
\index{documentation}

Every public item requires documentation:

\begin{lstlisting}[caption={Documentation example}, label=lst:docs]
/// A geometric point in 2D space.
///
/// Points are the fundamental building blocks of geometric constructions.
/// They can be created directly or as the result of operations like
/// line intersections.
///
/// # Mathematical Background
///
/// In our type-theoretic interpretation, a point corresponds to a
/// term of type `Point`, which is a dependent pair:
///
/// ```text
/// Point := Σ(x : ℝ)(y : ℝ)
/// ```
///
/// # Accessibility
///
/// Points are announced to screen readers with their coordinates
/// and any semantic labels (e.g., "intersection point A").
///
/// # Examples
///
/// ```rust
/// use proveit::geometry::Point;
///
/// let origin = Point::new(0.0, 0.0);
/// let p = Point::new(3.0, 4.0);
///
/// assert_eq!(origin.distance_to(&p), 5.0);
/// ```
#[derive(Clone, Copy, Debug, PartialEq)]
pub struct Point {
    /// The x-coordinate
    pub x: f64,
    /// The y-coordinate
    pub y: f64,
}
\end{lstlisting}

%-------------------------------------------------------------------------------
\chapter{Testing}
\label{chap:testing}
%-------------------------------------------------------------------------------

\section{Test Organization}
\index{testing}

\begin{lstlisting}[caption={Test module structure}, label=lst:tests]
#[cfg(test)]
mod tests {
    use super::*;
    use proptest::prelude::*;

    #[test]
    fn test_point_construction() {
        let p = Point::new(1.0, 2.0);
        assert_eq!(p.x, 1.0);
        assert_eq!(p.y, 2.0);
    }

    // Property-based tests for mathematical invariants
    proptest! {
        #[test]
        fn distance_is_symmetric(
            x1 in -1000.0..1000.0f64,
            y1 in -1000.0..1000.0f64,
            x2 in -1000.0..1000.0f64,
            y2 in -1000.0..1000.0f64,
        ) {
            let p1 = Point::new(x1, y1);
            let p2 = Point::new(x2, y2);

            let d1 = p1.distance_to(&p2);
            let d2 = p2.distance_to(&p1);

            prop_assert!((d1 - d2).abs() < 1e-10);
        }

        #[test]
        fn triangle_inequality(
            x1 in -100.0..100.0f64,
            y1 in -100.0..100.0f64,
            x2 in -100.0..100.0f64,
            y2 in -100.0..100.0f64,
            x3 in -100.0..100.0f64,
            y3 in -100.0..100.0f64,
        ) {
            let a = Point::new(x1, y1);
            let b = Point::new(x2, y2);
            let c = Point::new(x3, y3);

            let ab = a.distance_to(&b);
            let bc = b.distance_to(&c);
            let ac = a.distance_to(&c);

            prop_assert!(ac <= ab + bc + 1e-10);
        }
    }
}
\end{lstlisting}

%-------------------------------------------------------------------------------
\chapter{Appendices}
\label{chap:appendices}
%-------------------------------------------------------------------------------

\appendix

\chapter{Symbol Reference}
\label{app:symbols}

\begin{longtable}{@{}clp{8cm}@{}}
\caption{Mathematical symbols used in this manual} \\
\toprule
\textbf{Symbol} & \textbf{Name} & \textbf{Meaning} \\
\midrule
\endfirsthead
\multicolumn{3}{c}{\tablename\ \thetable{} -- continued} \\
\toprule
\textbf{Symbol} & \textbf{Name} & \textbf{Meaning} \\
\midrule
\endhead
\bottomrule
\endfoot
$\Type$ & Type & The universe of types \\
$\Pi$ & Pi (product) & Dependent function type \\
$\Sigma$ & Sigma (sum) & Dependent pair type \\
$\lambda$ & Lambda & Function abstraction \\
$\equiv$ & Definitional equality & Judgmental/computational equality \\
$=$ & Propositional equality & Identity type \\
$\simeq$ & Equivalence & Type equivalence \\
$\proves$ & Turnstile & Derivability \\
$\Hom{\mathcal{C}}{A}{B}$ & Hom-set & Morphisms from $A$ to $B$ \\
$\circ$ & Composition & Morphism composition \\
\end{longtable}

\chapter{Building from Source}
\label{app:building}

\section{Prerequisites}

\begin{itemize}
    \item Rust 1.70 or later
    \item Cargo package manager
    \item (Optional) \LaTeX{} distribution for building documentation
\end{itemize}

\section{Build Commands}

\begin{lstlisting}[language=bash, caption={Build commands}]
# Clone the repository
git clone https://github.com/TensorHusker/ProveIt.git
cd ProveIt

# Build in debug mode
cargo build

# Build in release mode
cargo build --release

# Run tests
cargo test

# Run with logging
RUST_LOG=debug cargo run

# Build documentation
cargo doc --open

# Build this manual (requires XeLaTeX)
cd docs/tex
xelatex proveit-manual.tex
xelatex proveit-manual.tex  # Run twice for cross-references
\end{lstlisting}

%-------------------------------------------------------------------------------
% INDEX
%-------------------------------------------------------------------------------
\printindex

\end{document}
